\documentclass[12pt]{amsart} 

\input{./latex-preamble/cm-preamble.tex}

%\graphicspath{ {./diagrams/} }


\newtheorem{lemma}{Lemma}

\begin{document}

{Hartshorne} \hfill {3-08-2025}

\bigskip\bigskip

\begin{center}

{\sc Chapter II.5, Week 7}

\end{center}

\begin{center}

  {Noah Parker} %\footnote{Collaborated with Kalev Martinson and Frank :).}}
    
\end{center}

\bigskip\bigskip

\begin{enumerate}
  \item[5.9.] {
      Let $S$ be a graded ring, generated by $S_1$ as an $S_0$-algebra, $M$ a graded $S$-module, and $X=\Proj S$.
      \begin{enumerate}
        \item Show there is a natural homomorphism $\alpha:M\to \Gamma_*(\wtilde{M})$.
        \item Assume now that $S_0=A$ is a finitely generated $k$-algebra for some field $k$, that $S_1$ is a finitely generated $A$-module, and that $M$ is a finitely generated $S$-module.
          Show that the map $\alpha$ is an isomorphism in all large enough degrees, i.e., there is a $d_0\in \ZZ$ such that for all $d\geqs d_0$, $\alpha_d:M_d\to \Gamma(X,\wtilde{M}(d))$ is an isomorphism.
          [\emph{Hint}: Use the methods of the proof of (5.19).]
        \item With the same hypothesis, we define an equivalence relation $\approx$ on graded $S$-modules by saying $M\approx M'$ if there is an integer $d$ such that $M_{\geqs d}\cong M_{\geqs d}'$.
          Here, $M_{\geqs d}=\bigoplus_{n\geqs d}M_n$.
          We say that a graded $S$-module is \emph{quasi-finitely generated} if it is equivalent to a finitely generated module.
          Now show that the functors $\wtilde\cdot$ and $\Gamma_*$ induce an equivalence of categories between the category of quasi-finitely generated graded $S$-modules modulo the equivalence relation $\approx$, and the category of coherent $\O_X$-modules.
      \end{enumerate}
    }
\end{enumerate}

\begin{proof}
  (a) Fix generators $\{x_i\}_{i\in I}\subset S_1$ of $S$ as an $S_0$-algebra.
  This gives us an affine open cover $\{D_+(x_i)\cong \Spec S_{(x_i)}\}_{i\in I}$ of $\Proj S$ such that 
  \begin{align*}
    \wtilde{M}|_{D_+(x_i)}(n)
    &\cong \wtilde{M(n)}|_{D_+(x_i)} \\
    &\cong (M(n)_{(x_i)})^\sim \\
    &\cong ((M_{x_i})_n)^\sim
  \end{align*}
  since restriction commute with tensor products, and $x_i\in S_1$ implies 
  \[
    (M\otimes_S N)_{(x_i)}=M_{(x_i)}\otimes_{S_{(x_i)}}N_{(x_i)}.
  \]
  Then
  \begin{align*}
    \Gamma_*(\wtilde{M}|_{D_+(x_i)})
    &\cong \bigoplus_{n\in \ZZ} \Gamma(D_+(x_i), \wtilde{M}|_{D_+(x_i)}(n)) \\
    %&\cong \bigoplus_{n\in \ZZ} \Gamma(D_+(x_i), \wtilde{M(n)}|_{D_+(x_i)}) \\
    &\cong \bigoplus_{n\in \ZZ} \Gamma(\Spec S_{(x_i)}, ((M_{x_i})_n)^\sim) \\
    &\cong \bigoplus_{n\in \ZZ}(M_{x_i})_n \\
    &\cong M_{x_i}
  \end{align*}
  supplies us with a natural map $\alpha_i:M\hookrightarrow M_{x_i}\iso \Gamma_*(\wtilde{M}|_{D_+(x_i)})$ for each $i$.
  Using a similar argument to the above, we define maps $\alpha_{ij}:M\hookrightarrow M_{x_ix_j}\iso \Gamma_+(\wtilde{M}|_{D_+(x_ix_j)})$ and show that $\rho_{D_+(x_ix_j)}\circ \alpha_i= \alpha_{ij} =\rho_{D_+(x_ix_j)} \circ \alpha_j$.

  Fixing $m\in M$, the $\alpha_i(m)_n\in \Gamma(D_+(x_i),\wtilde{M}|_{D_+(x_i}(n))$ glue together to a section $\alpha(m)_n\in \Gamma(X,\wtilde{M})$.
  Only finitely many $\alpha_i(m)_n$ are nonzero for each $i$, so only finitely many $\alpha(m)_n$ are nonzero, whence $\alpha(m)=\sum_n\alpha(m)_n\in \Gamma_*(\wtilde{M})$ is well-defined.
  We take $\alpha:m\mapsto \alpha(m)$ to be our natural map.

  (b) By (I, 7.4), we have a finite filtration
  \[
    0=M^0\subset M^1\subset\cdots\subset M^r=M
  \] 
  of $M$ by graded submodules, where for each $i$, $M^i/M^{i-1}\cong (S/\p_i)(n_i)$ for some homogeneous prime ideals $\p_i\lideal S$ and some integer $n_i$.
  Combining this with exactness of $\wtilde{\cdot}$, left exactness of $\Gamma(X,\cdot)$, and exactness of taking degree gives us the following diagram of $S_0$-modules with exact rows, for any $d\in\ZZ$.
  \[
    \begin{tikzcd}
      0\ar[r] & M^{i-1}_d\ar[r]\ar[d] & M^i_d\ar[r]\ar[d] &M^i/M^{i-1}_d\ar[r]\ar[d] & 0 \\
      0\ar[r] & \Gamma(X,\wtilde M^{i-1}(d))\ar[r]  &\Gamma(X,\wtilde M^i(d))\ar[r] & \Gamma(X,\wtilde{M^i/M^{i-1}}(d))
    \end{tikzcd}
  \] 
  If we have that the left and right vertical maps are isomorphisms, the five lemma implies the middle map is an isomorphism.

  The result clearly holds for $M^0=0$, so it suffices to show the result for for the $M^i/M^{i-1}\cong (S/\p_i)(n)$.
  Our map at degree $d$ is an isomorphism iff it is a bijection, allowing us to disregard our original $S$ and reduce to the case of $M=S$ a graded integral domain, finitely generated by $S_1$ as an $S_0$-algebra, and $S_0=A$ a finitely generated integral domain over $k$.

  Take such $S$, with $x_0,\ldots, x_r\in S_1$ generators as an $A$-module.
  Define $T_{\geqs n}:=\bigoplus_{d\geqs n}T_n$ for any graded ring $T$, and denote $S':=\Gamma_*(\O_X)_{\geqs 0}$.
  We will show that 
  \[
    \alpha_{\geqs m}:S_{\geqs m}\to S'_{\geqs m}
  \]
  is an isomorphism for $m\gg0$.
  %A section $s\in \Gamma(X,\O_X(d))$ consists of $s_i\in S_{x_i}$ for all $i$ such that $s_i=s_j\in S_{x_ix_j}$.
  $S$ is an integral domain, so the localization maps $S\to S_f$ are injective for all $f\neq 0$.
  In particular, each of the $S\to S_{x_i}$, $S_{x_i}\to S_{x_ix_j}$ are injective, as is $S\to S_{x_0\cdots x_r}$.
  Thus, $\alpha_{d\geqs 0}:S\to S'$ corresponds to the inclusion $S\subset S'\subset \cap_iS_{x_i}$.
  It suffices to show $S'_{\geqs m}\subset S_{\geqs m}$ for $m\gg0$.

  By the proof of 5.19, we have that $S'$ is a finitely generated $S$-module.
  Take such a set of homogeneous generators $y_0,\ldots, y_l$ of degree $d_j\geqs 0$.
  Since $y_j\in S_{x_i}$ for all $i$, we can find a single integer $n$ such that $x_i^ny_j\in S$ for all $i$.
  The $x_i$ generate $S_1$, so the monomials in the $x_i$ of degree $m$ generate $S_m$ for any $m$.
  Then we can take larger $n$ such that $fy_j\in S$ for all $f\in S_n$.
  Since $y_j$ is of positive degree, $fy_j\in S_{\geqs n}:=\bigoplus_{d\geqs n}S_d$ whenever $f\in S_{\geqs n}$.
  %It follows inductively that $fy_j^q\in S_{\geqs n}$ whenever $f\in S_{\geqs n}$.
  Take one $n$ that works for all $y_j$, so that $S_{\geqs n}\subset S_{\geqs n}[y_0,\ldots, y_l]\subset S_{\geqs n}$.
  Then take $m=n+\max_j d_j$, so that
  \begin{align*}
    S'_{\geqs m} \subset S_{\geqs n}[y_0,\ldots,y_l]=S_{\geqs n}\subset S_{\geqs m}.
  \end{align*}
  We conclude that $\alpha_{\geqs m}:S_{\geqs m}\to S'_{\geqs m}$ is an isomorphism for $m\gg0$, as desired.

  (c) If $M_d=0$ for $d\gg 0$, then any $m\in \Gamma(X,\wtilde{M})$ restricts to $m_i\in M_{x_i}$ with $m=\frac{x_i^d m}{x_i^d}=0$, whence $\wtilde{M}=0$.
  Note that $\ker(\phi_d)=(\ker\phi)_d$ and $\cok(\phi_d)=(\ker\phi)_d$, so the following sequence is unambiguous and exact.
  \[
    \begin{tikzcd}
      \ker\phi_{\geqs d}\ar[r] & M_{\geqs d}\ar[r,"\phi_{\geqs d}"] & M'_{\geqs d}\ar[r] & \cok\phi_{\geqs d}
    \end{tikzcd}
  \] 
  Then if $\phi_{\geqs d}:M_{\geqs d}\to M'_{\geqs d}$ is an isomorphism for $d\gg 0$, exactness of $\wtilde{\cdot}$ gives us
  \[
    \begin{tikzcd}
      0=\wtilde{\ker\phi}_{\geqs d}\ar[r] & \wtilde M_{\geqs d}\ar[r,"\phi_{\geqs d}"] & \wtilde{M'}_{\geqs d}\ar[r] & \wtilde{\cok\phi}_{\geqs d}=0
    \end{tikzcd}
  \] 
  is exact, whence $\wtilde M\cong \wtilde{M'}$.
  Thus, $\wtilde \cdot$ induces a map from $\approx$-classes to quasi-coherent $\O_X$-modules.
  Under the assumption that $S$ is noetherian, $\wtilde \cdot$ induces a map from $\approx$-classes of quasi-finitely generated $S$-modules to coherent $\O_X$-modules.

  We can use 5.19 to show that $\Gamma_*(\cdot)$ induces a map from coherent $\O_X$-modules to $\approx$-classes of quasi-finitely generated $S$-modules as well.
  By 5.15, we have a natural isomorphism $\beta:\Gamma_*(\cdot)^\sim\to \id$, and by (b) we have a natural isomorphism of $\approx$-classes $\alpha:\id\to \Gamma_*(\wtilde{\cdot})$.
  We conclude that $\Gamma_*(\cdot)$ and $\wtilde{\cdot}$ induce an equivalence of categories, as desired.
\end{proof}

\begin{enumerate}
  \item[5.10.] {
      Let $A$ be a ring, $S=A[x_0,\ldots, x_r]$, and $X=\Proj S$.
      We have seen that a homogeneous ideal $I$ in $S$ defines a closed subscheme of $X$ (Ex. 3.12), and that convrsely every closed subscheme of $X$ arises in this way (5.16).
      \begin{enumerate}
        \item For any homogeneous ideal $I\subset S$, we define the saturation $\ol{I}$ of $I$ to be
          \[
            \ol{I}:=\{s\in S\mid \forall 0\leqs i\leqs r~(\exists n>0~(x_i^ns\in I))\}
          \] 
          We say that $I$ is \emph{saturated} if $I=\ol{I}$.
          Show that $\ol{I}$ is a homogeneous ideal of $S$.
        \item Two homogeneous ideals $I_1$ and $I_2$ of $S$ define the same closed subscheme of $X$ iff they have the same saturation.
        \item If $Y$ is any closed subscheme of $X$, then the ideal $\Gamma_*(\I_Y)$ is saturated.
          Hence, it is the largest homogeneous ideal defining the subscheme $Y$.
        \item There is a 1-1 correspondence between saturated ideals of $S$ and closed sub-schemes of $X$.
      \end{enumerate}
    }
\end{enumerate}

\begin{proof}
  (a) If $s\in \ol{I}$, then we can write $s= \sum_ds_d$ as the sum of its homogeneous parts, so that $x_i^ns=\sum_dx_i^ns_d\in I$ and $I$ homogeneous implies $x_i^ns_d\in I$ for each $d$.
  We conclude that $s_d\in \ol{I}$ for all $d$, whence $\ol{I}$ is homogeneous.

  (b) Fix homogeneous ideals $I_1,I_2\lideal S$. 
  We have an affine cover $D_+(x_i)\cong \Spec S_{(x_i)}$ of $X$, so it suffices to show that 
  \[
    \frac{S_{(x_i)}}{(I_1)_{(x_i)}}
    \cong \wtilde{S/I_1}|_{D_+(x_i)}
    \cong \wtilde{S/I_2}|_{D_+(x_i)} 
    \cong \frac{S_{(x_i)}}{(I_2)_{(x_i)}}
  \] 
  for all $0\leqs i\leqs r$ iff $\ol{I}_1=\ol I_2$.
  This occurs iff the kernels of the respective projection maps are the same, i.e. $(I_1)_{(x_i)}= (I_2)_{(x_i)}$.
  Clearly this holds if $(I_1)_{x_i}=(I_2)_{x_i}$, so suppose conversely that the equality holds only in degree 0.
  Then $x_i$ is a unit, so
  \begin{align*}
    f\in ((I_1)_{x_i})_d 
    &\iff f/x_i^d\in (I_1)_{x_i}=(I_2)_{x_i} \\
    &\iff f\in((I_2)_{x_i})_d
  \end{align*}
  gives us that $(I_1)_{x_i}=(I_2)_{x_i}$.

  We have that
  \begin{align*}
    S\cap (I_1)_{x_i} &= \{s\in S:\exists n>0~ (x_i^ns\in I_1)\},
  \end{align*}
  so 
  \[
    \ol I_1 = S\cap \bigcap_{i=0}^r(I_1)_{x_i} =S\cap\bigcap_{i=0}^r(I_2)_{x_i}=\ol I_2
  \] 
  whenever the equality of localizations holds.
  Conversely, suppose $(I_1)_{x_i}\neq (I_2)_{x_i}$ so, without loss of generality, there exists $s\in (I_1)_{x_i}\setminus (I_2)_{x_i}$.
  Then we can take $n>0$ such that $x_i^ns\in I_1\subset\ol I_1$, but $x_i^m(x_i^ns)\notin I_2$ for all $m>0$ implies $x_i^ns\in \ol I_1\setminus \ol I_2$, whence $\ol I_1\neq \ol I_2$.

  (c) We have that $\I_Y=\ker (i^\sharp:\O_Y\to i_*\O_X)$ is a subsheaf of $\O_Y$, so 
  \begin{align*}
    J_i:=\Gamma(D_+(x_i),\I_Y)&\subset \Gamma(D_+(x_i),\O_Y)=S_{(x_i)}
  \end{align*}
  implies
  \begin{align*}
    J_i(n)\cong \Gamma(D_+(x_i),\I_Y(n))&\subset \Gamma(D_+(x_i),\O_Y(n))=S(n)_{(x_i)}
  \end{align*}
  for all $n\in \ZZ$,
  whence $\Gamma_*(\I_Y)\subset \Gamma_*(\O_X)\cong S$ by 5.15.

  Now, fix $s\in \ol J_m$, so that $x_i^rs\in J$ for $r\gg0$.
  Write
  \[
    x_i^mJ_i=J_i(m)\subset S(m)_{(x_i)}=x_i^m S_{(x_i)}.
  \] 
  We have $s_i:=s|_{D_+(x_i)}\in x_i^mS(m)_{(x_i)}$, and $x_i^rs_i \in (\Gamma_*(\I_Y)|_{D_+(x_i)})_{m+r}=J_{m+r}$.
  Then $x_i^rs_i=x_i^{m+r}j_i$ for $j_i\in J_i$, so $x_i\in J_i^\times$ gives us
  \[
    s_i=x_i^rj_i\in J_m.
  \] 
  Then $s$ is a gluing of sections in $s_i\in\Gamma(D_+(x_i),\I_Y(m))$, so $s\in \Gamma(X,\I_Y(m))\subset \Gamma_*(\I_Y)$, as desired.

  (d) To any saturated ideal $I\lideal S$, we associated the closed subscheme $Y_I:=\Proj S/I$.
  By (b), this map $I\mapsto Y_I$ is injective.
  By (c), it is surjective, with $\Gamma_*(\I_Y)\mapsto \Proj S/\Gamma_*(\I_Y)\cong Y$ for any closed subscheme $Y$.
  This gives us our bijection.
  
\end{proof}

% \begin{enumerate}
%   \item[5.11.] {
%       Let $S$ and $T$ be two graded rings with $S_0=T_0=A$.
%       We define the \emph{Cartesian product} $S\times_A T$ to be the graded ring $\bigoplus_{d\geqs 0}S_d\otimes_A T_d.$
%       If $X=\Proj S$ and $Y=\Proj T$, show that $\Proj (S\times_A T) \cong X\times_A Y$, and show that the sheaf $\O(1)$ on $\Proj(S\times_A T)$ is isomorphic to the sheaf $p_1^*(\O_X(1))\otimes p_2^*(\O_Y(1))$ on $X\times_A Y$.
% 
%       The Cartesian product of rings is related to the Segre embedding of projective spaces in the following way.
%       If $x_0,\ldots, x_r$ is a set of generated for $S_1$ over $A$, corresponding to a projective embedding $X\hookrightarrow \PP_A^r$, and if $y_0,\ldots, y_s$ is a set of generators for $T_1$, corresponding to a projective embedding $Y\hookrightarrow \PP_A^s$, then $x_i\otimes y_j$ is a set of generators for $(S\times_A T)_1$, and hence defines a projective embedding $\Proj (S\times_A T)\hookrightarrow \PP_A^N$, with $N=rs + r + s-1$.
%       This is just the image of $X\times Y\subset \PP^r\times \PP^s$ in its Sege embedding.
%     }
% \end{enumerate}
% 
% \begin{proof}
%   Take generators $\{x_i\}_{i\in I}\subset S_+$ and $\{y_j\}_{j\in J}\subset T_+$ of $S$ and $T$ as $A$-algebras.
%   Then any element $s\in S_d$ can be written $s$
% \end{proof}
% 
% \begin{enumerate}
%   \item[5.12.] {
%       \begin{enumerate}
%         \item Let $X$ be a scheme over a scheme $Y$, and let $\L$, $\M$ be two very ample invertible sheaves on $X$.
%           Show that $\L\otimes \M$ is also very ample.
%           [\emph{Hint}: Use a Segre embedding.]
%         \item Let $f:X\to Y$ and $g:Y\to Z$ be two morphisms of schemes.
%           Let $\L$ be a very ample invertible sheaf on $X$ relative to $Y$, and let $\M$ be a very ample invertible sheaf on $Y$ relative to $Z$.
%           Show that $\L\otimes f^*\M$ is a very ample invertible sheaf on $X$ relative to $Z$.
%       \end{enumerate}
%     }
% \end{enumerate}
% 
% \begin{proof}
%   (a) If $\L$, $\M$ are very ample, then we have $i:X\to \PP_Y^m$ and $j:X\to \PP_Y^n$ such that
% \end{proof}
% 
% \begin{enumerate}
%   \item[5.13.] {
%       Let $S$ be a graded ring, generated by $S_1$ as an $S_0$-algebra.
%       For any integer $d>0$, let $S^{(d)}$ be the graded ring $\bigoplus_{n\geqs 0}S_n^{(d)}$, where $S_n^{(d)}=S_{nd}$.
%       Let $X=\Proj S$.
%       Show that $\Proj S^{(d)}\cong X$, and that the sheaf $\O(1)$ on $\Proj S^{(d)}$ corresponds via this isomorphism to $\O_X(d)$.
%       
%       This construction is related to the $d$-uple \emph{embedding} in the following way.
%       Denote $A=S_0$.
%       If $x_0,\ldots, x_r$ is a set of generators for $S_1$, corresponding to an embedding $X\hookrightarrow \PP_A^r$, then the set of monomials of degree $d$ in the $x_i$ is a set of generators for $S_1^{(d)}=D_d$.
%       These define a projective embedding of $\Proj S^{(d)}$ which is none other than the image of $X$ under the $d$-uple embedding of $\PP_A^r$.
%     }
% \end{enumerate}
% 
% \begin{enumerate}
%   \item[5.14.] {
%       Let $A$ be a ring, and let $X$ be a closed subscheme of $\PP_A^r$.
%       We define the homogeneous coordinate ring $S(X)$ of $X$ for the given embedding to be $A[x_0,\ldots, x_r]/I$, where $I$ is the ideal $\Gamma_*(\I_X)$ constructed in the proof of (5.14).
%       (Of course, if $A$ is a field and $X$ a variety, this coincides with the usual definition!)
%       Recall that a scheme $X$ is \emph{normal} if its local rings are integrall closed domains.
%       A closed subscheme $X\subset \PP_A^r$ is \emph{projectively normal} for the given embedding, if the homogeneous coordinate ring $S(X)$ is an integrally closed domain .
%       Now assume that $A$ is a finitely generated $k$-algebra for some field $k$, and that $X$ is a connected, normal closed susbcheme of $\PP_A^r$.
%       Show that for some $d>0$, the $d$-uple embedding of $X$ is projectively normal, as follows.
%       \begin{enumerate}
%         \item Let $S$ be the homogeneous coordinate ring of $X$, and let $S'=\bigoplus_{n\geqs 0}\Gamma(X,\O_X(n))$.
%           Show that $S$ is a domain, and that $S'$ is its integral closure.
%           [\emph{Hint}: First show that $X$ is integral.
%           Then regard $S'$ as the global sections of the sheaf of rings $\sS=\bigoplus_{n\geqs 0}\O_X(n)$ on $X$, and show that $\sS$ is a sheaf of integrally closed domains.]
%         \item Use (Ex. 5.9) to show that $S_d=S_d'$ for all sufficiently large $d$.
%         \item Show that $S^{(d)}$ is integrally closed for sufficiently large $d$, and hence conclude that the $d$-uple embedding of $X$ is projectively normal.
%         \item As a corolary of (a), show that a closed subscheme $X\subset \PP_A^r$ is projectively normal if and only if it is normal, and for every $n\geqs 0$ the natural map $\Gamma(\PP^r,\O_{\PP^r}(n))\to \Gamma(X,\O_X(n))$ is surjective.a
%       \end{enumerate}
%     }
% \end{enumerate}

\end{document}
